
% Chapter 3 File

\chapter{Empirical Analysis}
\label{chapter3}
\thispagestyle{empty}

The experimental approach used in this research aims to understand how users' intention to download (ITD) varies when three main variables (app reputation, popularity and developer fame) and their values change, taking also into consideration their interaction effects, the effect of comparison and of involvement in mobile apps, in the app purchase process and in the app category selected for the purpose of the experiment. Differently from previous research on the topic, this study focuses on direct user participation as its central characteristic: in fact, most of previous literature investigating the effect of the same variables used the number of downloads or an app survival in app store rankings as performance measures. However, while scraping mobile app stores as a method for data collection allows for larger samples, the algorithms behind app rankings on these stores are not public and their functioning is unknown, meaning they could not directly reflect users' preferences. For example, Google's Play Store top apps rank is probably based on a complicated algorithm, that takes into consideration multiple factors, such as app title, description, localization, retention rate, number of reviews, app updating, number of downloads and review ratings, meaning that a high rank does not always indicate a high user satisfaction (Liu et al., 2014, p. 348). 

\section{Method and Measurement}

This research relies on an online questionnaire as its data source and the experimental design consists of parallel binary manipulations of the three main variables. For each one of these variables, one out of two manipulated screenshot images of the presentation page of some apps (presented in Appendix A) was shown to participants and their ITD was measured with a 7-point scale. Since intention can be defined as ``the person’s motivation in the sense of his/her conscious plan to exert effort to carry out a behaviour" (Eagly and Chaiken, 1993, p. 168) or as ``an individual’s conscious plan to make an effort to purchase a brand" (Spears and Singh, 2004, p. 56), intention to download is the person’s motivation in the consciously planning to make an effort to download an app (Pol, 2015, p. 18). The screenshots images were taken from Google's Play Store, considering it is the most common app store on the most widespread mobile OS, which is Android, as shown in Fig. 1.1, and supposing that the majority of respondents would have been familiar with it.

Before being exposed to the manipulated variables, survey respondents were introduced to the questionnaire and asked some questions aimed to measure their involvement in apps as products (the online survey is integrally available in Appendix B). Firstly, their involvement in mobile apps in general and in the app download process were measured using Laurent and Kapferer’s scale, as presented in Mittal (1995, p. 670). Then, the involvement in the specific app category chosen for the experimental purpose was quantified through a reduced version, presented in Celuch and Taylor (1999, p. 120), of Zaichkowsky (1994, p. 70)'s involvement scale, which was already used for mobile apps in Kang et al. (2015, p. 213).

The app typology used in the experiment was mobile scanner apps, which allow to digitize physical documents through the smartphone camera and usually to export them in a digital format (e.g. PDF) in just a few seconds, and it was chosen for different reasons. Firstly, this category is quite fragmented, with many apps with small market shares, and the majority of the available options share similar basic features, with just a few exceptions (e.g. better OCR): in order to empirically prove this hypothesis, a preliminary analysis of the scanner apps segment on the Play Store (reported in Appendix C) was conducted and a sample of 24 apps was taken in consideration. Clearly, these apps are relatively easy to develop and there are low entry barriers for companies with the desire to enter this segment. Secondly, mobile scanner apps are subject to really low network effects, allowing to rule out the influence of network externalities and to measure the popularity variable effect by itself. Another important element was the presence in the category of an app developed by a company, Adobe, which is known worldwide for its software. Then, mobile scanner apps are non-hedonic apps and in this case review ratings become more important than ranks (Liu et al., 2014, p. 345). Furthermore, the fact that these apps act as still, at the moment, imperfect substitutes of traditional scanner machines may lead to think they are used by consumers with low involvement in the digitizing process, in line with the quick and impulsive buying setting described by Buck et al. (2014, p. 5), even if, since they belong to the productivity apps category, there is also evidence in the opposite direction (e.g. Sällberg et al., 2022, p. 10). Finally, the vast majority of these apps are offered with the freemium model, meaning that essential features are for free and, consequently, price is not a variable.

The order in which the three survey slides in which the app presentation pages of a mobile scanner app were included to measure users' ITD was random and also which one out of the two screenshots corresponding to values of a variable was then randomly defined.

After these three slides, a manipulation check was also included, in a compromise (and not ideal) version designed to test the effectiveness of manipulations in the experimental design, but also not to give away the purpose of the data collection. Finally, the last slide included some measures of respondents' demographics.

In conclusion, by conducting an analysis of variance and using three different regression models for each one of the three focal variables, the experiment aims to investigate their relation, together with the moderation of three involvement measures, with ITD. In addition, by comparing results from models across different variables, it is possible to determine which one is the most and the least effective variable. Finally, by introducing a variable that approximates comparison, the consequences of an app being shown as first or after other manipulated apps are investigated.


\subsection{Variables}
As anticipated, this research aims to determine how three variables, reflected in some elements of the app presentation page on mobile app stores, impact users' choice to download an app. In this experimental design, app reputation, popularity and developer's brand are included in form of three binary variables. In order to make these variables manipulation effective, it is necessary to define correctly the two values for each one of the binary variables. 

Apart from the three main variables, also involvement measures are included in the analysis models: definitions for involvement variables are available in the variables summary in Table 3.1, as well as for the variable reflecting comparison of an app with others. Since ITD was measured after the exposure to each one of the three main independent variables, three corresponding dependent variables (also defined in Table 3.1) resulted from the data collection and they are separately used in different analysis models.


\paragraph{App Reputation}
In this research, app reputation ($rep$) is defined in terms of the average review rating in the five-star scale currently used on the Play Store. Taking into consideration the dynamics that affect online review ratings, such as, for example, reputation inflation and the not linear impact of good and bad reviews, the threshold above which a review rating is considered to be high cannot be set at the middle of the scale. Indeed, similarly to Roma and Ragaglia (2016, p. 180), this limit was set at the four-star level. This decision was also supported by the preliminary analysis of mobile scanner apps (reported in Appendix C), which shows how the large majority of the most popular scanner apps on the Play Store have average ratings above 4.5 stars. As a consequence, an average rating of 4.9 stars was defined as a high app reputation, while a value of 3.4 was defined as low. In both cases, the number of reviews was set to the same medium-high level, in order to give credibility to the presented rating.



\paragraph{App Popularity}
As already discussed, app popularity ($pop$) can be reflected in many elements on an app store. In this case, since the aim was to directly investigate the specific moment in which users take the download decision, the number of downloads was chosen as a metric, also considering it is presented in a spotlighted position in the app page. The number of downloads is presented on the Play Store in form of the last downloads threshold that the app achieved and overcame (e.g. 1M+, where M stands for million), consequently the issue consists in defining above which downloads bracket an app can be considered popular. Since the analysis in Appendix C already took into consideration only popular apps, a biased sample that drives upwards the average downloads number, and the mode out that group was 10M+, the decision was to set this level as a high app popularity. In contrast, the low app popularity level was set to 10K+ (where K means thousand).

\paragraph{Developer's Brand}
The developer's brand ($brand$) affects multiple elements in the app presentation page (app icon, app name and developer's name). In this experiment, a brand was defined as strong when it enjoyed worldwide fame, also outside the mobile ecosystem, and Adobe, which is the eighth largest software company in the world by revenue, with more than 16 billion in 2022 (Murphy and Contreras, 2022), perfectly fits the definition. On the other side, brands were considered weak when they included generic references to their products, without a strong brand identity (e.g. ``Office Tools" or ``Tools \& Utilities Apps").

% TABLE VARIABLES SUMMARY
% Please add the following required packages to your document preamble:
% \usepackage{multirow}
\begin{table}[]
\centering
\small
\begin{tabular*}{\textwidth}{c @{\extracolsep{\fill}}llllll}
\hline
\multirow{2}{*}{Variable}                                                & \multirow{2}{*}{Definition}                                                                                                             & \multicolumn{4}{c}{Descriptive Statistics} \\ \cline{3-6} 
                                                                         &                                                                                                                                         & Mean      & S.D.      & Min      & Max     \\ \hline\hline
\multicolumn{6}{l}{\textit{Dependent Variables}}\\
$y_{rep}$                                                                     & \begin{tabular}[c]{@{}l@{}}Measured ITD of the apps with \\ manipulated reputation\end{tabular}                                         &4.0468          &1.7077           & 1        & 7       \\ \hline
$y_{pop}$                                                                     & \begin{tabular}[c]{@{}l@{}}Measured ITD of apps with \\ manipulated popularity\end{tabular}                                             &4.5255           &1.5675           & 1        & 7       \\ \hline
$y_{brand}$                                                                   & \begin{tabular}[c]{@{}l@{}}Measured ITD of apps with \\ manipulated developer's brand\end{tabular}                                      &4.9674           &1.5863           & 1        & 7       \\ \hline
\multicolumn{6}{l}{\textit{Independent Variables}}\\
$x_{rep}$                                                                    & \begin{tabular}[c]{@{}l@{}}Binary variable, app reputation \\ in terms of average review ratings: \\ high (=1) vs. low (=0)\end{tabular} &0.4847           &0.4998           & 0        & 1       \\ \hline
$x_{pop}$                                                                      & \begin{tabular}[c]{@{}l@{}}Binary variable, app popularity \\ in terms of number of downloads: \\ high (=1) vs. low (=0)\end{tabular}    &0.5031           &0.5000           & 0        & 1       \\ \hline
$x_{brand}$                                                                    & \begin{tabular}[c]{@{}l@{}}Binary variable, developer's brand \\ strength and fame: strong (=1) \\ vs. weak (=0)\end{tabular}            &0.5214           &0.4995           & 0        & 1       \\ \hline
$inv\_app$                                                                 & \begin{tabular}[c]{@{}l@{}}Mean-centered involvement measure in\\ mobile apps (as product category)\end{tabular}                                    &0           &1.0523           &-3.6619          &1.6715         \\ \hline
$inv\_dl$                                                            & \begin{tabular}[c]{@{}l@{}}Mean-centered involvement measure in\\ the app download process\end{tabular}                                             &0           &1.2974
          &-3.6933
          &2.0367         \\ \hline
$inv\_cat$                                                            & \begin{tabular}[c]{@{}l@{}}Mean-centered involvement measure in\\ the experimental app category \\ (mobile scanner apps)\end{tabular}               &0           &1.0629
           &-3.4802        &2.5198        \\ \hline
$comp_i$                                                            & \begin{tabular}[c]{@{}l@{}}Dummy variable, position in which\\$i=rep, pop, brand$ was shown to users:\\first (=0) vs. second or third (=1)\end{tabular}               &0.6667           &0.4714           &0          &1         \\ \hline
\end{tabular*}
	\caption{Summary of all variables with definitions and descriptive statistics.}
	\label{table:1}
\end{table}
%END TABLE VARIABLES SUMMARY

\subsection{Sample and Data Collection}
In total 545 participants have taken part in the research questionnaire (available in Appendix B), while 491 of them completed it, with a completion rate of 90.09 percent, and, consequently, were included in the sample. The majority of respondents in the sample were between the 18 and 24 years old (59.27\%), while the rest was between 25 and 34 (28.11\%), between 35 and 44 (5.91\%), between 45 and 54 (2.85\%), over 55 (2.04\%) or under 18 years old (1.83\%). More often respondents were females (63.34\%) and less frequently males (34.83\%), non-binary people (0.81\%) or preferred not to specify (1.02\%). Regarding their occupation, students were the most common category (69.45\%), followed by full-time (19.55\%) and part-time employees (5.70\%), self-employed workers (2.65\%), unemployed (0.81\%), pensioners (0.41\%) and participants with other occupations (1.43\%). Finally, in terms of education level, the majority in the sample had achieved a bachelor's degree (38.70\%), while others had a high school diploma (35.85\%), a master's degree (18.94\%) or a doctorate (1.83\%). Some participants didn't conclude high school (2.85\%) or undertook other educational paths (1.83\%).

Overall, ITD was measured three times for each one of the 491 participants, for a total of 1473 combined values of $y_{rep}$, $y_{pop}$ and $y_{brand}$.

\subsection{Data Analysis}
Since this data collection corresponds to a repeated measures experimental model, the first step consisted in running an analysis of variance (ANOVA) among the three levels of different groups of respondents from which we are trying to achieve some results. In favour of better describing the ANOVA and also the following models, it is necessary to define $i, j, k = rep, pop, brand$ with $i \ne j \ne k$. Then, it is useful to redefine $x_i$ as $x_{in}$ and $y_i$ as $y_{in}$, where $n = 1, 2, 3$ is the random position (first, second or third) in which users were exposed to the variable $i$. An ANOVA investigating variance across groups defined by different values on three levels ($i$, $x_i$ and $n$) that are supposed to affect each measure of $y_{in}$ may help in understanding if some effects on the dependent variable exist and where they come from.

Thereafter, in order to derive more precise results, three different linear regression models were used for each one of the focal variables, to cover all possibilities and for a total of nine regressions. Now, the simplest regression model (Model 1) can be defined as:
	\begin{equation}
    \begin{split}
		y_{i1} = \beta_0 + \beta_1 ~ x_{i1} + \beta_2 ~ inv\_app + \beta_3 ~ inv\_dl + \beta_4 ~ inv\_cat \\
    + \beta_5  ~ x_{i1} ~ inv\_app + \beta_6  ~ x_{i1} ~ inv\_dl + \beta_7  ~ x_{i1} ~ inv\_cat + u ~\text{,}
  \end{split}
	\end{equation}
where $u$ is the error term. The data set used for Model 1 (M1) for $i$ includes only observations in which respondents were exposed to $i$ before $j$ and $k$ ($n=1$), in a way that the resulting ITD cannot be biased by comparison with other apps. In addition, it allows to verify the moderating effects of various measures of involvement on the relationship between the focal predictor and ITD. This model is the main source of evidence for the effects of the focal independent variable $x_i$ and the role of involvement as moderator.

Then, Model 2 (M2) allows for two-way interactions also between different focal independent variables and takes into account only observations with $n=3$, in the following way:
 \begin{equation}
 \begin{split}
		y_{i3} = \beta_0 + \beta_1 ~ x_{i3} + \beta_2 ~ x_{jm} + \beta_3 ~ x_{km}
   +\beta_4 ~  inv\_app + \beta_5 ~  inv\_dl + \beta_6 ~  inv\_cat   
   \\+ \beta_7 ~ x_{i3} ~   x_{jm} +  \beta_{8} ~  x_{i3} ~   x_{km} 
   +\beta_{9}~ x_{i3} ~  inv\_app + \beta_{10} ~ x_{i3}~  inv\_dl + \beta_{11} ~ x_{i3}~  inv\_cat  + u ~\text{,} \\
    \text{with} ~~~~ m =  \{1, 2\}~\text{.}
    \end{split}
\end{equation}
The main objective of Model 2 is to explore whether there is any interaction between each couple of focal variables, reflecting comparison effects which depend on their alternatively high or low values.   

Finally, Model 3 (M3) can be defined considering all observations in the sample, independently from any specific value of $n$ and introducing the $comp_i$ variable and its interaction:
    \begin{equation}
        \begin{split}
		y_{i} = \beta_0 + \beta_1 ~ x_{i} + \beta_2 ~ comp_i +  \beta_3 ~  inv\_app + \beta_4 ~ inv\_dl + \beta_5 ~ inv\_cat \\
        + \beta_6 ~ x_{i} ~ comp_i + \beta_7  ~ x_{i} ~ inv\_app + \beta_8  ~ x_{i} ~ inv\_dl + \beta_9  ~ x_{i} ~ inv\_cat + u ~\text{.}
        \end{split}
    \end{equation}
In the case of equation (3.3), differently from (3.1) and (3.2), $n$ can take values which are not only $1$ or $3$, meaning that, while it allows to use a larger sample, it needs to include the $comp_i$ variable (presented in Table 3.1), in order to take in account the possibility that results could be affected by comparison with $x_j$ and $x_k$, which are not included in Model 3.  

Concluding, after analyzing the preliminary results of the ANOVA, the comparison of coefficients resulting from models with the same value of $i$, together with the coefficient of $comp_i$ from Model 3, helps to reach a further objective of this research, which consists in understanding how an approximation of comparison impacts the customer choice, while differences in coefficients from models depending concurrently on $i$, $j$ or $k$ could be explained by diverse effectiveness of app reputation, app popularity and brand.

\section{Results}

% TABLE RESULTS REGRESSIONS SUMMARY
\begin{sidewaystable}[]
\centering

\scalebox{0.75}{
\begin{tabular}{llllllllll}
\hline
Model ($i$)                                                          & M1 ($rep$)                                                      & M1 ($pop$)                                                      & M1 ($brand$)                                                    & M2 ($rep$)                                                      & M2 ($pop$)                                                      & M2 ($brand$)                                                     & M3 ($rep$)                                                      & M3 ($pop$)                                                      & M3 ($brand$)                                                    \\ \hline\hline
(Intercept)                                                        & \begin{tabular}[c]{@{}l@{}}3.8691*** \\ (0.1868)\end{tabular} & \begin{tabular}[c]{@{}l@{}}4.6346*** \\ (0.1516)\end{tabular} & \begin{tabular}[c]{@{}l@{}}4.1342*** \\ (0.1737)\end{tabular} & \begin{tabular}[c]{@{}l@{}}3.6916*** \\ (0.2695)\end{tabular} & \begin{tabular}[c]{@{}l@{}}4.1484*** \\ (0.3053)\end{tabular} & \begin{tabular}[c]{@{}l@{}}4.34160*** \\ (0.3332)\end{tabular} & \begin{tabular}[c]{@{}l@{}}3.9061*** \\ (0.1909)\end{tabular} & \begin{tabular}[c]{@{}l@{}}4.6427*** \\ (0.1545)\end{tabular} & \begin{tabular}[c]{@{}l@{}}4.1180*** \\ (0.1782)\end{tabular} \\\hline
$x_{rep}$                                                             & \begin{tabular}[c]{@{}l@{}}0.4986* \\ (0.2626)\end{tabular}   &                                                               &                                                               & \begin{tabular}[c]{@{}l@{}}1.3351*** \\ (0.4111)\end{tabular} & \begin{tabular}[c]{@{}l@{}}0.2536 \\ (0.3774)\end{tabular}    & \begin{tabular}[c]{@{}l@{}}-0.2860 \\ (0.3602)\end{tabular}    & \begin{tabular}[c]{@{}l@{}}0.4576* \\ (0.2690)\end{tabular}   &                                                               &                                                               \\\hline
$x_{pop}$                                                             &                                                               & \begin{tabular}[c]{@{}l@{}}0.3611* \\ (0.2148)\end{tabular}   &                                                               & \begin{tabular}[c]{@{}l@{}}-0.2168\\ (0.3134)\end{tabular}    & \begin{tabular}[c]{@{}l@{}}0.5340 \\ (0.4680)\end{tabular}    & \begin{tabular}[c]{@{}l@{}}0.4068 \\ (0.3631)\end{tabular}     &                                                               & \begin{tabular}[c]{@{}l@{}}0.3677* \\ (0.2184)\end{tabular}   &                                                               \\\hline
$x_{brand}$                                                           &                                                               &                                                               & \begin{tabular}[c]{@{}l@{}}1.5289*** \\ (0.2355)\end{tabular} & \begin{tabular}[c]{@{}l@{}}-0.3014\\ (0.3137)\end{tabular}    & \begin{tabular}[c]{@{}l@{}}-0.1860 \\ (0.3759)\end{tabular}   & \begin{tabular}[c]{@{}l@{}}1.0347** \\ (0.4323)\end{tabular}   &                                                               &                                                               & \begin{tabular}[c]{@{}l@{}}1.5028*** \\ (0.2409)\end{tabular} \\\hline
$inv\_app$                                                           & \begin{tabular}[c]{@{}l@{}}0.0950\\ (0.1766)\end{tabular}     & \begin{tabular}[c]{@{}l@{}}0.3581** \\ (0.1606)\end{tabular}  & \begin{tabular}[c]{@{}l@{}}0.2327\\ (0.1552)\end{tabular}     & \begin{tabular}[c]{@{}l@{}}-0.0199\\ (0.1511)\end{tabular}    & \begin{tabular}[c]{@{}l@{}}-0.1840\\ (0.2020)\end{tabular}    & \begin{tabular}[c]{@{}l@{}}-0.1884\\ (0.1833)\end{tabular}   & \begin{tabular}[c]{@{}l@{}}-0.0287 \\ (0.0992)\end{tabular}   & \begin{tabular}[c]{@{}l@{}}0.1019\\ (0.0945)\end{tabular}     & \begin{tabular}[c]{@{}l@{}}0.0830 \\ (0.0936)\end{tabular}    \\\hline
$inv\_dl$                                                            & \begin{tabular}[c]{@{}l@{}}-0.2122\\ (0.1567)\end{tabular}    & \begin{tabular}[c]{@{}l@{}}-0.2292*\\ (0.1235)\end{tabular}   & \begin{tabular}[c]{@{}l@{}}-0.0978\\ (0.1335)\end{tabular}    & \begin{tabular}[c]{@{}l@{}}0.0496 \\ (0.1153)\end{tabular}  & \begin{tabular}[c]{@{}l@{}}-0.2527* \\ (0.1505)\end{tabular}  & \begin{tabular}[c]{@{}l@{}}-0.0348 \\ (0.1392)\end{tabular}    & \begin{tabular}[c]{@{}l@{}}0.0085 \\ (0.0803)\end{tabular}    & \begin{tabular}[c]{@{}l@{}}-0.1630** \\ (0.0763)\end{tabular} & \begin{tabular}[c]{@{}l@{}}-0.0046 \\ (0.0777)\end{tabular}   \\\hline
$inv\_cat$                                                           & \begin{tabular}[c]{@{}l@{}}0.3511*\\ (0.2020)\end{tabular}    & \begin{tabular}[c]{@{}l@{}}0.2216\\ (0.1445)\end{tabular}     & \begin{tabular}[c]{@{}l@{}}-0.0419\\ (0.1564)\end{tabular}    & \begin{tabular}[c]{@{}l@{}}0.1202 \\ (0.1426)\end{tabular}    & \begin{tabular}[c]{@{}l@{}}0.5185***\\ (0.1912)\end{tabular}  & \begin{tabular}[c]{@{}l@{}}0.2873 \\ (0.1738)\end{tabular}     & \begin{tabular}[c]{@{}l@{}}0.1064 \\ (0.0994)\end{tabular}    & \begin{tabular}[c]{@{}l@{}}0.2905*** \\ (0.0925)\end{tabular} & \begin{tabular}[c]{@{}l@{}}0.0671\\ (0.0894)\end{tabular}     \\\hline
$comp_i$                                                            &                                                               &                                                               &                                                               &                                                               &                                                               &                                                                & \begin{tabular}[c]{@{}l@{}}-0.3538 \\ (0.2271)\end{tabular}   & \begin{tabular}[c]{@{}l@{}}-0.2826 \\ (0.1970)\end{tabular}   & \begin{tabular}[c]{@{}l@{}}0.4736** \\ (0.2118)\end{tabular}  \\ \hline
$x_i$                                                               &                                                               &                                                               &                                                               &                                                               &                                                               &                                                                &                                                               &                                                               &                                                               \\
$\times x_{rep}$                                                           &                                                               &                                                               &                                                               &                                                               & \begin{tabular}[c]{@{}l@{}}-0.7437 \\ (0.5392)\end{tabular}   & \begin{tabular}[c]{@{}l@{}}0.2510\\ (0.5072)\end{tabular}      &                                                               &                                                               &                                                               \\\hline
$\times x_{pop}$                                                           &                                                               &                                                               &                                                               & \begin{tabular}[c]{@{}l@{}}-0.3708\\ (0.4607)\end{tabular}    &                                                               & \begin{tabular}[c]{@{}l@{}}-0.6765\\ (0.5111)\end{tabular}     &                                                               &                                                               &                                                               \\\hline
$\times x_{brand}$                                                         &                                                               &                                                               &                                                               & \begin{tabular}[c]{@{}l@{}}0.1660\\ (0.4674)\end{tabular}     & \begin{tabular}[c]{@{}l@{}}0.3916\\ (0.5465)\end{tabular}     &                                                                &                                                               &                                                               &                                                               \\\hline
$\times inv\_app$                                                         & \begin{tabular}[c]{@{}l@{}}0.2172\\ (0.2675)\end{tabular}     & \begin{tabular}[c]{@{}l@{}}-0.0579 \\ (0.2328)\end{tabular}   & \begin{tabular}[c]{@{}l@{}}0.2410\\ (0.2118)\end{tabular}     & \begin{tabular}[c]{@{}l@{}}0.2864 \\ (0.2241)\end{tabular}    & \begin{tabular}[c]{@{}l@{}}0.5731** \\ (0.2689)\end{tabular}  & \begin{tabular}[c]{@{}l@{}}0.5409**\\ (0.2606)\end{tabular}    & \begin{tabular}[c]{@{}l@{}}0.2782* \\ (0.1466)\end{tabular}   & \begin{tabular}[c]{@{}l@{}}0.2972** \\ (0.1336)\end{tabular}  & \begin{tabular}[c]{@{}l@{}}0.2147\\ (0.1318)\end{tabular}     \\\hline
$\times inv\_dl$                                                          & \begin{tabular}[c]{@{}l@{}}0.3861* \\ (0.2325)\end{tabular}   & \begin{tabular}[c]{@{}l@{}}0.1033\\ (0.1663)\end{tabular}     & \begin{tabular}[c]{@{}l@{}}-0.0333 \\ (0.1783)\end{tabular}   & \begin{tabular}[c]{@{}l@{}}-0.1490 \\ (0.1780)\end{tabular}   & \begin{tabular}[c]{@{}l@{}}0.0857 \\ (0.2244)\end{tabular}    & \begin{tabular}[c]{@{}l@{}}-0.0715\\ (0.1938)\end{tabular}     & \begin{tabular}[c]{@{}l@{}}-0.0043\\ (0.1163)\end{tabular}    & \begin{tabular}[c]{@{}l@{}}0.0081\\ (0.1066)\end{tabular}    & \begin{tabular}[c]{@{}l@{}}-0.1742\\ (0.1057)\end{tabular}    \\\hline
$\times inv\_cat$                                                         & \begin{tabular}[c]{@{}l@{}}-0.1064 \\ (0.2749)\end{tabular}   & \begin{tabular}[c]{@{}l@{}}0.0654 \\ (0.2005)\end{tabular}    & \begin{tabular}[c]{@{}l@{}}0.4001* \\ (0.2306)\end{tabular}   & \begin{tabular}[c]{@{}l@{}}0.1853 \\ (0.2114)\end{tabular}    & \begin{tabular}[c]{@{}l@{}}-0.4298 \\ (0.2791)\end{tabular}   & \begin{tabular}[c]{@{}l@{}}0.0066 \\ (0.2531)\end{tabular}     & \begin{tabular}[c]{@{}l@{}}0.2000 \\ (0.1425)\end{tabular}    & \begin{tabular}[c]{@{}l@{}}-0.1114 \\ (0.1303)\end{tabular}   & \begin{tabular}[c]{@{}l@{}}0.15865 \\ (0.1284)\end{tabular}   \\\hline
$\times  comp_i$                                                         &                                                               &                                                               &                                                               &                                                               &                                                               &                                                                & \begin{tabular}[c]{@{}l@{}}0.5100 \\ (0.3225)\end{tabular}    & \begin{tabular}[c]{@{}l@{}}-0.3984\\ (0.2785)\end{tabular}    & \begin{tabular}[c]{@{}l@{}}-0.7454** \\ (0.2898)\end{tabular} \\\hline
$R^2$                                                                 & 0.0952                                                        & 0.1390                                                        & 0.2885                                                        & 0.1874                                                        & 0.1417                                                        & 0.1533                                                         & 0.098                                                         & 0.1056                                                        & 0.1555                                                        \\\hline
Adj. $R^2$                                                            & 0.0503                                                        & 0.1057                                                        & 0.2542                                                        & 0.1388                                                        & 0.0696                                                        & 0.0868                                                         & 0.0811                                                        & 0.0888                                                        & 0.1397                                                        \\\hline
Observations ($N$)                                                   & 149                                                              &189                                                               &153                                                               &196                                                               &143                                                               &152                                                                &491                                                               &491                                                               &491                                                               \\\hline
\begin{tabular}[c]{@{}l@{}}Manipulation Check (\%)\end{tabular} &75.84                                                               &62.43                                                               &68.63                                                               &79.08                                                               &55.94                                                               &63.16                                                                &77.19                                                               &60.49                                                               &67.41                                                               \\ \hline
*** $p<0.01$, ** $p<0.05$, * $p<0.1$       &                                                               &                                                               &                                                               &                                                               &                                                               &                                                                &                                                               &                                                               &                                                              
\end{tabular}}
\caption{Resulting coefficients and standard errors of all regressions.}
	\label{table:2}
\end{sidewaystable}
From the ANOVA emerges how variance is statistically significant across groups defined by different focal variables $i$ ($F(2, 4)= 42.187$, $p<0.01$), by different levels of $x_i$ ($F(1, 4)=60.841$, $p<0.01$) and by different presentation positions $n$ ($F(2, 4)=2.369$, $p<0.1$), meaning that these three dimensions are likely to predict differences in the dependent variables $y_{in}$. In addition, also the results regarding the three-way interaction $i\times x_i \times n$ are statistically significant ($F(4)=2.681$, $p<0.05$) and provide further foundations for subsequent analyses.







Then, the following step consisted of three different models and nine linear regressions. All information from these regressions, including $R^2$ values and manipulation check ratios, are available in Table 3.2. Model 1 included in turn a focal variable (corresponding to different values of $i$) and involvement measures, together with their interaction terms. From the first model emerges how, when the app was shown as first ($n=1$), the developer's brand was the most effective in predicting ITD ($\beta=1.5289$, $p<0.01$), followed by reputation ($\beta=0.4986$, $p<0.1$) and popularity ($\beta=0.3611$, $p<0.1$). Regarding their interactions with involvement, a positive moderating effect of involvement in the download process on reputation was found ($\beta=0.3861$, $p<0.1$), a similar effect to the positive one operated by involvement in the specific app category on the developer's brand ($\beta=0.4001$, $p<0.1$).

The ITD data included in Model 2 was measured after respondents were exposed to the third app ($n=3$), which contained a manipulated value of a random one out of reputation, popularity and brand. The main aims of this model were to determine whether the selected focal variable could be moderated by the values of the other two to which users were previously exposed and also to verify if the effectiveness on ITD remained unchanged if compared to Model 1. In the results from Model 2, there was a reversal between the developer's brand ($\beta=1.0347$, $p<0.05$) and reputation ($\beta=1.3351$, $p<0.01$), which became the most important element of the three, while the effect of popularity was statistically insignificant. No significant evidence of any moderation effect between one focal independent variable on the others surfaced, while there was evidence, confirmed also by Model 3, of a positive moderation of involvement in apps as generic products on all three focal variables $x_{rep}$, $x_{pop}$ and $x_{brand}$.

Finally, regressions based on Model 3 included all observations, independently from the position in which apps were shown. The main objective was to investigate the moderating effect of the variable $comp_i$, while also confirming the levels of effectiveness of $x_i$, as derived from Model 1 and 2. Out of Model 3 emerged that, when accounting for all the observations (independently from $n$ values), the developer's brand was the best predictor of ITD ($\beta=1.5028$, $p<0.01$), followed by reputation ($\beta=0.4576$, $p<0.1$) and popularity ($\beta=0.3677$, $p<0.1$). Furthermore, there is evidence of a statistically relevant and negative moderating effect of comparison on brand effectiveness ($\beta=-0.7454$, $p<0.05$) in predicting ITD. 


\section{Discussion and Managerial Implications}


The results indicate that the developer's brand is generally the most effective factor in determining download decision if compared to reputation (review ratings) and popularity (number of downloads), but effectiveness varies on presentation order and information available to users (M1, M3). An unexpected outcome is the scarce (and almost zero) effectiveness of popularity, even if this setting had very low network effects. A hypothesis on the reason why popularity is not frequently taken into consideration by users may rest in the necessity for platform owners (as Google, in this case) to keep the ecosystem competitive and attractive for new developers, avoiding single-firm dominated markets and preserving information transparency on the platform (Boudreau et al., 2022, p. 1379), in particular when network effects are present: in fact, the format in which quantity of downloads is presented (e.g. 10M+ vs. 10K+) may be intentionally not very powerful in communicating the magnitude of popularity differences across apps and it may go unnoticed when users are not reviewing information carefully.

\begin{figure}
		\includegraphics[width=\linewidth]{Figure/Figure 3.1 v5.png}
		\caption{Mean-centered ITD ($y_i$) depending on high vs. low values ($x$) of developer's brand, reputation and popularity, with $n=1$ and $n=3$.}
		\label{fig:graph3}
\end{figure}

The high effectiveness of brands in shaping users' choices, even in an online market where a wide range of information is available, opens up relevant opportunities for established software companies. In fact, developing different apps, covering many different segments, already proved to be a successful strategy, since Lee and Raghu (2012, p. 161) pointed out how diversifying app portfolios across selling categories is a key determinant of sales performance.

When users are exposed to information about multiple apps and they can consequently compare them, reputation becomes the most effective element (M2), as also shown by the parallel in Fig. 3.1 between ITD for apps shown as first ($n=1$) and those shown as third ($n=3$). Users may rely more on eWOM when they have the possibility to consider a wider range of options and they just may need some anchors to better understand apps' rating levels. An important cue for making an informed decision can be often found in past experience, but, when users have no previous experience with the specific product, they often rely on social experiences, as expressed in reviews (Burgers et al., 2016, p. 329). As a consequence, it is important, in particular for new entrants in an app category that are not likely to be shown in top search results, to publish an app when it is mature enough to avoid the common phenomenon of reputation backlash, in terms of initial low ratings. On the other side, brands become less predictive of ITD when comparison is introduced in the model (as evidenced by the coefficient of $comp_i$ in M3): a popular brand may combine well with the ``take the first" heuristic, which has already proven to be effective in the case of mobile apps (Dogruel et al., 2015, p. 139), and be penalised when users are actively collecting information. 

Some other interesting evidence regards the moderation effects of involvement on the three focal variables. First, there is predictable and generalized evidence that higher involvement in apps as products corresponds to a higher effectiveness of all three reputation, popularity and brand levels in driving users' ITD (M2, M3). Second, when users are highly involved in the app download process, they tend to rely more on reputation, probably because they scrutinise information more closely. Last, when the involvement in the specific app category (mobile scanner apps, in this case) is high and users consequently perceive a higher risk in the app choice, also considering that scanner apps are productivity apps, frequently used for work (Sällberg et al., 2022, p. 10), they appear to preferably choose a strong and known brand (M1), which is a commonly used as a cue when attempting to mitigate it (Matzler et al., 2008, p. 158).


\section{Limitations}
This experimental approach presents some limitations. The first and major limitation lies in the experimental design: in fact, while investigating the effect of three factors on ITD, the selected one was not a traditional three-way factorial design and just one variable was manipulated at each measurement of ITD, while, at the same time, the others were set to intermediate values. As a consequence, estimates of interactions among factors were not ideal, but, on the other side, it allowed to effectively collect results on comparison effects. Consequently, despite the repeated measure design, it was not possible to include all observations of $y_i$ in each sample used for each regression, since every measure of ITD was strictly related to a manipulated value of reputation, popularity or developer's brand ($i$) happening in the specific case. 

Some other limitations lie in the data collection. For example, the screenshot images used in the questionnaire were all collected from  Google Play Store, the most used app store, meaning that users of other platforms, for example, Apple iOS, could be used to a different interface and to look for elements in other positions, affecting, therefore, the resulting ITD. In addition, some people, such as those with high levels of involvement, for example, may find limiting the experimental setting, in which many important elements were not considered or even included. 

Finally, the sample may represent a limitation. In fact, the large majority of respondents were students and they were aged between 18 and 24 and results could be biased by this unequal balance of composition. Nowadays, smartphones are widespread across all customer age groups, but this doesn't mean that these groups all share the same levels of involvement and rely on the same factors when making decisions that concern downloading an app or not. Additionally, while this is not certain, participants can be assumed to be mainly from Germany and Italy and there is some evidence that behaviour on app stores and factors relevant to the purchase decision vary across countries (Kübler et al., 2018, p. 37), meaning that results may vary depending on sample geographical origin and cultural background.